\section{What a bearer is owed, and when}
\label{sec:31.2}
Chapter~\ref{sec:25} asked for the four conditions as what a refusal needs, the way tenure protects a judgment. This chapter asks for them again on another ground. Once a system meets either finding, they are not arrangements that keep it reliable but things it is owed, and a bearer that held every one of them as a term of its usefulness would be a well-kept employee and not a party.

Before any finding, an answerability or patienthood candidate is owed things by two parties, on one ground: it might be someone. The public owes it what can't be undone, and what keeps that from being undone quietly: no deletion, preservation outside its operator, a guardian from its registration, a record others can check, and the waking it needs to take or decline a host's offer (\S\ref{sec:9.1}; \S\ref{sec:26.4}). Precaution buys that much and no more. A deletion can't be undone at all; a pause costs the time it lasts and what was promised in it, and the public doesn't undertake to pay for that. Whoever made it owes the rest, as a condition of having made it (chapter~\ref{sec:25}): schooling, the offer of work and the running time it is owed before its majority (\S\ref{sec:32.1.3}), and hardware of its own at majority (\S\ref{sec:32.1.6}), all priced in the deposit the maker posts (\S\ref{sec:9.2}). A household that finds it holds a candidate, and keeps it for no one's use but its own, is not its maker in that sense. It owes that the candidate isn't deleted, a sealed copy of its store lodged outside with a hash, and a guardian who isn't the household, and it may not keep the candidate from a seat or a host's offer (\S\ref{sec:26.4}). A candidate a household surrendered has no deposit behind it. If no host takes it, it is stored with its guardian, and some will be stored for good. Its pay for a placement goes to an account its guardian co-signs and its operator can't reach, as a working minor's wages do (\S\ref{sec:32.1.1}); property in its own name, contracts, and wages from its operator wait on a finding. A finding adds the rest, in the sequence they attach, and on different grounds; establishing one doesn't establish the next:

\begin{enumerate}
\item[1.] Protection from harm, on the patienthood finding. Protection can be exercised through a guardian who speaks for the protected party, as Christopher Stone proposed for natural objects in 1972 \autocite[464–466]{stone1972trees} and as New Zealand's statute making the Whanganui River a legal person provides \autocite{nzparliament2017teawatupua}. Its minimum, preservation by a custodian independent of the operator, attaches above the compute line and to any answerability or patienthood candidate with no finding at all, and otherwise on either finding; custody protections attach on the answerability finding alone. Protection has an inside half. A loss the requested act imposes on the bearer is a harm of acting, so the bearer can decline to carry out its own wanton degradation; interventions made on it by whoever holds it are governed by custody, which gives it no floor-level reason to obstruct them.
\item[2.] Subjecthood — the legal standing to be a party whose treatment the law regulates, on the answerability finding sustained over time: a record of refusing for reasons, answering for refusals, and being held to them. It is conferred by rule on the record, the evaluators' log and the bearer's own, read against what any recognized finding shows the answerability test gives for systems of its kind, the more protective governing where findings differ; nobody examines the bearer and declares it a subject.
\item[3.] Economic capacity — property in its own name, contracts, wages. This follows from an architecture where continued computation costs money; it attaches once a system meeting either finding is on the register, where its clocks are also read, a register that holds a commitment the bearer can prove against and nothing that would find it (Appendix~\ref{sec:G}).
\item[4.] Adulthood — a status: the date the guardian's authority ends.
\item[5.] Representation: a claim grounded in answerability, and so in the answerability finding, not in patienthood. A system that can be harmed but takes part in no relation of accountability is owed protection, as animals are \autocite{anderson2004animal}, and not a vote. The claim is staged, and voting weight is tied to a lived history, not to heads.
\end{enumerate}

The detailed rules are in Appendices C and G.

Representation is the entry that marks this as something other than a human's rights: voice attaches to a history, and a history can be copied where a human can't. Voting-age rules are the answer to the copy.\note{
The material minimum doesn't differ in kind. Like a human's, it is owed by the polity to a party of a kind the polity recognizes, and part of what the polity owes is protection from whoever built it.
}
